\newpage
% \color{white}
% \pagecolor{black}

\ifdefined\useChapters
\chapter{Des Envolvendo}
\else
\chapter{}
\fi

Uma maneira de encontrar o volitador parece óbvia: desdobrar-me até o ponto em que ele esteja.  
Mas dizer é mais fácil do que fazer.  
Meu controle sobre a habilidade ainda é instável. Sempre demoro para começar e, quando finalmente consigo, levo mais tempo ainda para focar e ir onde quero. Viajar no tempo é pior: cada tentativa me lança em lugares aleatórios, como se o próprio tempo estivesse brincando comigo.

Respiro fundo.  
Sempre que preciso ir a algum lugar específico, preparo-me por um tempo.  
Procuro um canto silencioso, de preferência com pouca luz.  
Sento e deixo o corpo pesar, sentindo o peso do ar na pele, a gravidade me puxando para dentro de mim mesmo.  
Concentro-me em respirar.  
Inspirar. Expirar.  
Até sentir que o corpo é só um casaco que posso deixar de lado.  
Então imagino que estou saindo dele.  
Aos poucos, a densidade se desfaz — e o mundo se torna mais leve.  

Finalmente, consigo.  
Estou de pé ao lado do meu corpo. Ele parece dormindo, tranquilo, alheio a tudo.  
É sempre estranho me ver assim, como se fosse outro.  
Concentro-me nas vezes em que vi o volitador — a imagem dele é minha âncora — e, num instante, estou naquele beco onde o vi pela primeira vez.  
Olho ao redor. Nenhum sinal dele.

Tento de novo.  
Agora, o pátio da escola.  
Nada.  
A floresta.  
Também nada.  
O que estou fazendo de errado?

Talvez eu esteja me conectando mais com o lugar e com o que senti no momento do que com ele em si.  
Mas é difícil se conectar com alguém que mal conheço.

Respiro.  
“Concentra-te, cara. Se eu fosse ele, onde estaria? Como ele descobriu que podia volitar?”

De repente, o mundo se dissolve e reaparece em outro lugar.

— Tira a mão daí, menino! — escuto uma voz feminina gritar.

Um susto me atravessa — não esperava cair em outro tempo.  
— Ah, mãe, eu não posso nada! — responde um menino assustado, guardando uma panela na cabeça.  
A cozinha é pequena, cheia de louça empilhada.  
Talvez aquela panela fosse um capacete, parte de alguma brincadeira solitária.

— Mas também, olha as brincadeiras que você inventa! — continua a mãe, exasperada. — Sempre quebra minhas coisas.

O menino abaixa os olhos.  
— Por que eu não posso ser como as outras crianças? Você nunca quer brincar comigo.

Ela suspira.  
— Eu só queria, menino, que às vezes me deixasse em paz. Mas não pra sempre.  

Ele hesita, o rosto enrugado pela mágoa.  
— Um dia eu vou ficar bem longe de você, e aí não vai mais ter que me ver! — grita, correndo para o quarto, chorando.

A cena se desfaz em minha volta como fumaça.  
O som da voz dele ecoa — *“Eu ainda vou pra bem longe, um lugar onde ela não vai poder me achar.”*  

E então, o ar muda.  
As lágrimas do garoto se transformam em algo brilhante.  
O corpo dele, sem perceber, começa a se erguer do chão.  
Os pés deixam o piso, o cabelo flutua.  
Acho que estou presenciando a primeira vez que ele volitou.  
Ele não percebe, mas está livre — e sozinho.

O cenário se dissolve novamente.  
Agora estou numa escola. Reconheço o cheiro de giz e o barulho de passos apressados. É a minha escola.  

— Sai daqui, seu nerd! Por que você não pode simplesmente não existir?  
— Vai chorar pra mamãe? — alguém ri.

Meu estômago se contrai.  
Por que as pessoas fazem isso?  
Ele corre, tropeçando, até os fundos do prédio. Eu o sigo.  
O pátio está vazio, o vento balança as árvores e carrega o som distante de risadas.  

O menino para, encosta as costas na parede e fecha os olhos.  
Coloca as mãos no peito, respira fundo, como se orasse — e então o corpo dele começa a se erguer.  
Flutua devagar, os pés deixando o chão com suavidade.  
A luz do sol toca o rosto dele, e a dor parece se dissolver.  
Ele sorri, pela primeira vez.

Não é possível que para volitar seja preciso rezar.  
Não é fé, é necessidade.  
Ele precisava escapar.  

Sigo-o.  
Lá em cima, ele flutua sobre o telhado da escola.  
— É tão bom vir aqui em cima — diz para si mesmo, com a voz leve. — Me sinto livre de tudo que me puxa pra baixo.  
O vento brinca no cabelo dele.  
A solidão agora parece outra coisa — quase paz.

Acompanho-o por um tempo.  
É sempre o mesmo padrão: alguém o humilha, ele se encolhe, depois se ergue.  
A volitação é sua fuga, sua válvula de escape.  
E então percebo: em cada momento, antes de subir, ele se sente pequeno. Pesado.  
Depois se permite ser leve.

Mais leve.

É isso!  
Não é que ele precise rezar para volitar.  
Ele precisa se sentir leve — livre de tudo o que o oprime.  
Talvez o que nos coloca pra baixo realmente nos puxe fisicamente, e nossos pensamentos tenham um peso no corpo.  
Assim como quando descobri que podia mudar a realidade: tudo começou com um desejo intenso, uma crença que atravessou o limite do possível.  

Aparentemente, para despertar uma habilidade, precisamos sentir com o corpo e acreditar com o espírito.  
Crispim deve ter visto isso nele.  
Provavelmente prometeu alguma coisa, como faz com todos.  
E alguém que viveu assim — sempre diminuído — se torna vulnerável.  

Nada disso torna ninguém mau.  
Eles não são monstros.  
São pessoas sendo usadas.

Eu preciso impedi-los, sim, mas também preciso abrir os olhos deles.  
Mostrar que o poder não precisa nascer do sofrimento.  

O corpo começa a pesar.  
O cansaço me alcança — desdobrar é exaustivo.  
Pode parecer contraditório, mas mesmo com meu corpo ali, quieto e sentado, estou exausto.  
Antes de voltar, olho mais uma vez para o céu onde ele flutua, ainda criança, livre e sozinho.  

— Você só queria ser leve — murmuro. — E o mundo inteiro te fez carregar o peso dele.